\section{Results}
\label{sec:results}

\subsection{Matched-budget comparison across models, criteria and sparsities}

Table~\ref{tab:main} is the paper's main result. Each cell pair reports test
accuracy for the I.P.\ baseline and for BaCP under an identical budget, together
with their difference $\Delta = \mathrm{BaCP} - \mathrm{I.P.}$ Rows are grouped by
pruning criterion and ordered by sparsity within each group; columns are grouped
by model. Because the two arms within a cell differ only in the training
objective (Section~\ref{sec:experiments}), $\Delta$ is the quantity this paper
measures, and the absolute columns are context for it.

\begin{table}[t]
\centering
\footnotesize
\setlength{\tabcolsep}{3.5pt}
\caption{Test accuracy (\%) on the held-out CIFAR-10 test set under a matched
budget. $\Delta = \mathrm{BaCP} - \mathrm{I.P.}$; positive favours BaCP. Sparsity
is the fraction of prunable weights masked; the classifier head is in the
prunable scope and the projection head is not. Cells are single-seed except the
headline cells of Table~\ref{tab:seeds}.}
\label{tab:main}
\begin{tabular}{lccc ccc ccc ccc}
\toprule
& \multicolumn{3}{c}{ResNet-34} & \multicolumn{3}{c}{ResNet-50}
& \multicolumn{3}{c}{VGG-11} & \multicolumn{3}{c}{VGG-19} \\
\cmidrule(lr){2-4}\cmidrule(lr){5-7}\cmidrule(lr){8-10}\cmidrule(lr){11-13}
Sparsity & I.P. & BaCP & $\Delta$ & I.P. & BaCP & $\Delta$
         & I.P. & BaCP & $\Delta$ & I.P. & BaCP & $\Delta$ \\
\midrule
\multicolumn{13}{l}{\textit{Dense reference (unpruned, cross-entropy)}} \\
--- & \multicolumn{3}{c}{\result{resnet34.dense.none.0.0}}
    & \multicolumn{3}{c}{\result{resnet50.dense.none.0.0}}
    & \multicolumn{3}{c}{\result{vgg11.dense.none.0.0}}
    & \multicolumn{3}{c}{\result{vgg19.dense.none.0.0}} \\
\midrule
\multicolumn{13}{l}{\textit{Magnitude} \citep{han2015learning}} \\
0.95  & \result{resnet34.ip.magnitude.0.95}  & \result{resnet34.bacp.magnitude.0.95}  & \result{resnet34.delta.magnitude.0.95}
      & \result{resnet50.ip.magnitude.0.95}  & \result{resnet50.bacp.magnitude.0.95}  & \result{resnet50.delta.magnitude.0.95}
      & \result{vgg11.ip.magnitude.0.95}     & \result{vgg11.bacp.magnitude.0.95}     & \result{vgg11.delta.magnitude.0.95}
      & \result{vgg19.ip.magnitude.0.95}     & \result{vgg19.bacp.magnitude.0.95}     & \result{vgg19.delta.magnitude.0.95} \\
0.97  & \result{resnet34.ip.magnitude.0.97}  & \result{resnet34.bacp.magnitude.0.97}  & \result{resnet34.delta.magnitude.0.97}
      & \result{resnet50.ip.magnitude.0.97}  & \result{resnet50.bacp.magnitude.0.97}  & \result{resnet50.delta.magnitude.0.97}
      & \result{vgg11.ip.magnitude.0.97}     & \result{vgg11.bacp.magnitude.0.97}     & \result{vgg11.delta.magnitude.0.97}
      & \result{vgg19.ip.magnitude.0.97}     & \result{vgg19.bacp.magnitude.0.97}     & \result{vgg19.delta.magnitude.0.97} \\
0.99  & \result{resnet34.ip.magnitude.0.99}  & \result{resnet34.bacp.magnitude.0.99}  & \result{resnet34.delta.magnitude.0.99}
      & \result{resnet50.ip.magnitude.0.99}  & \result{resnet50.bacp.magnitude.0.99}  & \result{resnet50.delta.magnitude.0.99}
      & \result{vgg11.ip.magnitude.0.99}     & \result{vgg11.bacp.magnitude.0.99}     & \result{vgg11.delta.magnitude.0.99}
      & \result{vgg19.ip.magnitude.0.99}     & \result{vgg19.bacp.magnitude.0.99}     & \result{vgg19.delta.magnitude.0.99} \\
0.999 & \result{resnet34.ip.magnitude.0.999} & \result{resnet34.bacp.magnitude.0.999} & \result{resnet34.delta.magnitude.0.999}
      & \result{resnet50.ip.magnitude.0.999} & \result{resnet50.bacp.magnitude.0.999} & \result{resnet50.delta.magnitude.0.999}
      & \result{vgg11.ip.magnitude.0.999}    & \result{vgg11.bacp.magnitude.0.999}    & \result{vgg11.delta.magnitude.0.999}
      & \result{vgg19.ip.magnitude.0.999}    & \result{vgg19.bacp.magnitude.0.999}    & \result{vgg19.delta.magnitude.0.999} \\
\midrule
\multicolumn{13}{l}{\textit{SNIP-it} \citep{lee2019snip,verdenius2020snipit}} \\
0.95  & \result{resnet34.ip.snip.0.95}  & \result{resnet34.bacp.snip.0.95}  & \result{resnet34.delta.snip.0.95}
      & \result{resnet50.ip.snip.0.95}  & \result{resnet50.bacp.snip.0.95}  & \result{resnet50.delta.snip.0.95}
      & \result{vgg11.ip.snip.0.95}     & \result{vgg11.bacp.snip.0.95}     & \result{vgg11.delta.snip.0.95}
      & \result{vgg19.ip.snip.0.95}     & \result{vgg19.bacp.snip.0.95}     & \result{vgg19.delta.snip.0.95} \\
0.97  & \result{resnet34.ip.snip.0.97}  & \result{resnet34.bacp.snip.0.97}  & \result{resnet34.delta.snip.0.97}
      & \result{resnet50.ip.snip.0.97}  & \result{resnet50.bacp.snip.0.97}  & \result{resnet50.delta.snip.0.97}
      & \result{vgg11.ip.snip.0.97}     & \result{vgg11.bacp.snip.0.97}     & \result{vgg11.delta.snip.0.97}
      & \result{vgg19.ip.snip.0.97}     & \result{vgg19.bacp.snip.0.97}     & \result{vgg19.delta.snip.0.97} \\
0.99  & \result{resnet34.ip.snip.0.99}  & \result{resnet34.bacp.snip.0.99}  & \result{resnet34.delta.snip.0.99}
      & \result{resnet50.ip.snip.0.99}  & \result{resnet50.bacp.snip.0.99}  & \result{resnet50.delta.snip.0.99}
      & \result{vgg11.ip.snip.0.99}     & \result{vgg11.bacp.snip.0.99}     & \result{vgg11.delta.snip.0.99}
      & \result{vgg19.ip.snip.0.99}     & \result{vgg19.bacp.snip.0.99}     & \result{vgg19.delta.snip.0.99} \\
0.999 & \result{resnet34.ip.snip.0.999} & \result{resnet34.bacp.snip.0.999} & \result{resnet34.delta.snip.0.999}
      & \result{resnet50.ip.snip.0.999} & \result{resnet50.bacp.snip.0.999} & \result{resnet50.delta.snip.0.999}
      & \result{vgg11.ip.snip.0.999}    & \result{vgg11.bacp.snip.0.999}    & \result{vgg11.delta.snip.0.999}
      & \result{vgg19.ip.snip.0.999}    & \result{vgg19.bacp.snip.0.999}    & \result{vgg19.delta.snip.0.999} \\
\midrule
\multicolumn{13}{l}{\textit{WANDA} \citep{sun2024wanda}, global threshold} \\
0.95  & \result{resnet34.ip.wanda.0.95}  & \result{resnet34.bacp.wanda.0.95}  & \result{resnet34.delta.wanda.0.95}
      & \result{resnet50.ip.wanda.0.95}  & \result{resnet50.bacp.wanda.0.95}  & \result{resnet50.delta.wanda.0.95}
      & \result{vgg11.ip.wanda.0.95}     & \result{vgg11.bacp.wanda.0.95}     & \result{vgg11.delta.wanda.0.95}
      & \result{vgg19.ip.wanda.0.95}     & \result{vgg19.bacp.wanda.0.95}     & \result{vgg19.delta.wanda.0.95} \\
0.97  & \result{resnet34.ip.wanda.0.97}  & \result{resnet34.bacp.wanda.0.97}  & \result{resnet34.delta.wanda.0.97}
      & \result{resnet50.ip.wanda.0.97}  & \result{resnet50.bacp.wanda.0.97}  & \result{resnet50.delta.wanda.0.97}
      & \result{vgg11.ip.wanda.0.97}     & \result{vgg11.bacp.wanda.0.97}     & \result{vgg11.delta.wanda.0.97}
      & \result{vgg19.ip.wanda.0.97}     & \result{vgg19.bacp.wanda.0.97}     & \result{vgg19.delta.wanda.0.97} \\
0.99  & \result{resnet34.ip.wanda.0.99}  & \result{resnet34.bacp.wanda.0.99}  & \result{resnet34.delta.wanda.0.99}
      & \result{resnet50.ip.wanda.0.99}  & \result{resnet50.bacp.wanda.0.99}  & \result{resnet50.delta.wanda.0.99}
      & \result{vgg11.ip.wanda.0.99}     & \result{vgg11.bacp.wanda.0.99}     & \result{vgg11.delta.wanda.0.99}
      & \result{vgg19.ip.wanda.0.99}     & \result{vgg19.bacp.wanda.0.99}     & \result{vgg19.delta.wanda.0.99} \\
0.999 & \result{resnet34.ip.wanda.0.999} & \result{resnet34.bacp.wanda.0.999} & \result{resnet34.delta.wanda.0.999}
      & \result{resnet50.ip.wanda.0.999} & \result{resnet50.bacp.wanda.0.999} & \result{resnet50.delta.wanda.0.999}
      & \result{vgg11.ip.wanda.0.999}    & \result{vgg11.bacp.wanda.0.999}    & \result{vgg11.delta.wanda.0.999}
      & \result{vgg19.ip.wanda.0.999}    & \result{vgg19.bacp.wanda.0.999}    & \result{vgg19.delta.wanda.0.999} \\
\bottomrule
\end{tabular}
\end{table}

We read Table~\ref{tab:main} along three axes.

\paragraph{Across sparsity.}
The comparison at 0.95 is the least informative one in the table, because both
arms are close to the dense reference and there is little damage for either
objective to repair. The interesting region is the right end of the schedule,
where the surviving weight budget is small enough that the retraining signal has
to do real work. The controlled measurement we already have
(Section~\ref{sec:results:ablation}) shows the margin growing as sparsity rises,
and this is the shape a representation-level effect predicts: the more of the
representation the mask destroys, the more there is for an objective that targets
the representation directly to recover. Table~\ref{tab:main} is the test of
whether that shape holds across models and criteria.
% TODO: confirm direction once the sweep lands -- the claim "the margin widened with sparsity" must be checked against every model column, not only ResNet-50.

\paragraph{Across criteria.}
The three criterion blocks share everything but the scoring function of
Eq.~\eqref{eq:criteria}, so a $\Delta$ that keeps its sign across all three blocks
is evidence that the contrastive objective is orthogonal to the choice of
criterion rather than an interaction with one particular saliency signal. This is
the second of our three claims and it is the one the table is structured to test:
consistency of sign across blocks is what matters here, not the ranking of the
criteria themselves.
% TODO: confirm direction once the sweep lands -- if any criterion block reverses sign, the criterion-agnostic claim must be narrowed to the criteria where it holds, and the reversal reported rather than averaged away.

\paragraph{Across architectures.}
The ResNet and VGG columns are not expected to behave identically. The VGGs carry
no normalisation layers and concentrate their parameters in the classifier MLP,
which is in the prunable scope; the ResNets are heavily normalised and
convolution-dominated. Where the two families diverge, the architectural
difference is the first explanation to check, and Section~\ref{sec:experiments:vgglr}
already documents one place where it required a protocol adjustment.
% TODO: confirm direction once the sweep lands -- whether the VGG and ResNet families agree in sign and in trend with sparsity.

\begin{table}[t]
\centering
\small
\caption{Headline cells, 3 seeds, mean $\pm$ standard deviation. These are the
only cells in the paper with seed replication; every other cell is single-seed.}
\label{tab:seeds}
\begin{tabular}{llccc}
\toprule
Model & Sparsity & I.P. & BaCP & $\Delta$ \\
\midrule
ResNet-50 & 0.95 & \resultpm{resnet50.ip.magnitude.0.95} & \resultpm{resnet50.bacp.magnitude.0.95} & \result{resnet50.delta.magnitude.0.95} \\
ResNet-50 & 0.99 & \resultpm{resnet50.ip.magnitude.0.99} & \resultpm{resnet50.bacp.magnitude.0.99} & \result{resnet50.delta.magnitude.0.99} \\
VGG-19    & 0.95 & \resultpm{vgg19.ip.magnitude.0.95}    & \resultpm{vgg19.bacp.magnitude.0.95}    & \result{vgg19.delta.magnitude.0.95} \\
VGG-19    & 0.99 & \resultpm{vgg19.ip.magnitude.0.99}    & \resultpm{vgg19.bacp.magnitude.0.99}    & \result{vgg19.delta.magnitude.0.99} \\
\bottomrule
\end{tabular}
\end{table}

Table~\ref{tab:seeds} reports the headline cells with seed replication. The
comparison that matters there is between the reported spread and the observed
$\Delta$: a difference that does not clear the seed spread is not a result, and we
report the spread so that a reader can apply that test directly rather than take
our word for it.
% TODO: confirm direction once the sweep lands -- if a headline delta does not clear its own seed spread, say so explicitly here.

\subsection{Which contrastive objective: a controlled ablation}
\label{sec:results:ablation}

Equation~\eqref{eq:functional} can be realised in more than one way, and the
choice is not neutral. Table~\ref{tab:ablation} reports a two-variable ablation on
ResNet-50 with the magnitude criterion, in which exactly two configuration keys
were changed --- the contrastive formulation and the projection-head mode --- and
everything else was held identical and verified programmatically before the runs.
The \emph{legacy} arm is the objective described in Section~\ref{sec:method}:
SupCon plus NT-Xent over the square $2B \times 2B$ matrix, with per-model
trainable projection heads. The \emph{control} arm is closer to
\citet{xu2022cap}: one contrastive functional over the rectangular $B \times N$
student--teacher matrix, with a single frozen projection head shared by student
and teachers so that the head cannot absorb the objective. Both are compared
against the shared I.P.\ baseline.

\begin{table}[t]
\centering
\small
\caption{Objective ablation. ResNet-50, CIFAR-10, magnitude criterion,
\textbf{single seed}, under an \textbf{earlier protocol} than
Table~\ref{tab:main} (60 pruning epochs, $\Delta T = 88$, no validation split,
50-epoch fine-tune). Reported for the ordering it establishes, not for its
absolute values. \emph{Control} is the CAP-form rectangular functional with a
single frozen projection head shared by student and teachers; \emph{Legacy} is
the objective of Section~\ref{sec:method}. The control arm at 0.99 was
terminated by an unrelated job cancellation and was not re-run.}
\label{tab:ablation}
\begin{tabular}{lcccc}
\toprule
Sparsity & I.P.\ baseline & Control & Legacy (ours) & Legacy $-$ I.P. \\
\midrule
0.95 & \result{resnet50.ablip.magnitude.0.95} & \result{resnet50.ablcontrol.magnitude.0.95} & \result{resnet50.abllegacy.magnitude.0.95} & \result{resnet50.abldelta.magnitude.0.95} \\
0.97 & \result{resnet50.ablip.magnitude.0.97} & \result{resnet50.ablcontrol.magnitude.0.97} & \result{resnet50.abllegacy.magnitude.0.97} & \result{resnet50.abldelta.magnitude.0.97} \\
0.99 & \result{resnet50.ablip.magnitude.0.99} & \result{resnet50.ablcontrol.magnitude.0.99} & \result{resnet50.abllegacy.magnitude.0.99} & \result{resnet50.abldelta.magnitude.0.99} \\
\bottomrule
\end{tabular}
\end{table}

Three things follow, and we state them at the strength the evidence supports.
First, the legacy formulation was ahead of the control at every measured
sparsity, which is why it is the objective this paper studies; the choice was made
by measurement, and we report the arm that lost rather than omitting it. Second,
both contrastive arms were at or above the I.P.\ baseline at every measured point,
so the direction of the effect does not depend on which realisation is used ---
only its size does. Third, the legacy arm's margin over I.P.\ grew as sparsity
rose, which is the trend a representation-preservation mechanism predicts and the
reason we extended the grid to 0.999.

The caveats are load-bearing. This ablation is \textbf{single-seed}, and while its
gap exceeds the measured noise floor of
\result{cifar10.noise.floor.lo}--\result{cifar10.noise.floor.hi} points it does not
exceed it by a wide margin. It was run under the \textbf{earlier protocol} listed
in the caption, so its absolute values are not comparable to Table~\ref{tab:main}
and only its internal ordering should be read. And the winning arm is the one with
the trainable projection head, which is precisely the configuration that
complicates the interpretation of the loss --- see Section~\ref{sec:discussion}.
The control arm is therefore not merely a losing variant; it is the variant under
which the representation-preservation reading of the loss would have been clean,
and it is reported here as the cost of that cleanliness.

\subsection{Calibration against published results}

Table~\ref{tab:published} places our cells beside published CIFAR-10 numbers for
the same architectures. Only \textbf{95\%} of our sparsity grid overlaps the
literature's, which evaluates at 90 / 95 / 98; at 97\%, 99\% and 99.9\% there is
no static published comparator for these models, and we mark those rows as
standalone. The published values are quoted exactly from their sources and are
\emph{not} matched-budget comparisons against our runs: as
Section~\ref{sec:experiments:budget} sets out, they were produced with roughly an
order of magnitude more optimizer updates than we spent.

\begin{table}[t]
\centering
\small
\caption{Our cells beside published CIFAR-10 results. Published values are quoted
from their sources under \emph{their} recipes (160 epochs, batch 128 for the
GraNet-table rows; 250 epochs, batch 128 for EAST), not ours. They calibrate our
baseline; they are not a matched comparison. GraNet and GraSP disagree on the
VGG-19 dense baseline and we report both rather than choosing.}
\label{tab:published}
\begin{tabular}{lllll}
\toprule
Model & Sparsity & Method & Accuracy (\%) & Source \\
\midrule
\multicolumn{5}{l}{\textit{ResNet-50 / CIFAR-10}} \\
ResNet-50 & dense & published dense       & $94.75 \pm 0.01$ & \citep{liu2021granet} \\
ResNet-50 & dense & ours                  & \result{resnet50.dense.none.0.0} & this paper \\
ResNet-50 & 0.95  & SNIP                  & $90.86$          & \citep{liu2021granet} \\
ResNet-50 & 0.95  & GraSP                 & $91.32$          & \citep{liu2021granet} \\
ResNet-50 & 0.95  & SynFlow               & $91.22$          & \citep{liu2021granet} \\
ResNet-50 & 0.95  & RigL (dynamic)        & $93.86 \pm 0.25$ & \citep{liu2021granet} \\
ResNet-50 & 0.95  & GraNet (dynamic)      & $94.44 \pm 0.01$ & \citep{liu2021granet} \\
ResNet-50 & 0.95  & GMP $=$ gradual magnitude & $94.52 \pm 0.08$ & \citep{liu2021granet} \\
ResNet-50 & 0.95  & \textbf{I.P.\ magnitude (ours)} & \result{resnet50.ip.magnitude.0.95} & this paper \\
ResNet-50 & 0.95  & \textbf{BaCP magnitude (ours)}  & \result{resnet50.bacp.magnitude.0.95} & this paper \\
\midrule
\multicolumn{5}{l}{\textit{VGG-19 / CIFAR-10}} \\
VGG-19 & dense & published dense (GraNet) & $93.85 \pm 0.05$ & \citep{liu2021granet} \\
VGG-19 & dense & published dense (GraSP)  & $94.23$          & \citep{wang2020grasp} \\
VGG-19 & dense & ours                     & \result{vgg19.dense.none.0.0} & this paper \\
VGG-19 & 0.90  & GMP                      & $93.59 \pm 0.10$ & \citep{liu2021granet} \\
VGG-19 & 0.90  & RigL (dynamic)           & $93.38 \pm 0.11$ & \citep{liu2021granet} \\
VGG-19 & 0.95  & SNIP                     & $93.43 \pm 0.20$ & \citep{liu2021granet,wang2020grasp} \\
VGG-19 & 0.95  & GraSP                    & $93.04 \pm 0.18$ & \citep{liu2021granet,wang2020grasp} \\
VGG-19 & 0.98  & SNIP                     & $92.05 \pm 0.28$ & \citep{liu2021granet,wang2020grasp} \\
VGG-19 & 0.98  & GraSP                    & $92.19 \pm 0.12$ & \citep{liu2021granet,wang2020grasp} \\
VGG-19 & 0.95  & \textbf{I.P.\ magnitude (ours)} & \result{vgg19.ip.magnitude.0.95} & this paper \\
VGG-19 & 0.95  & \textbf{BaCP magnitude (ours)}  & \result{vgg19.bacp.magnitude.0.95} & this paper \\
VGG-19 & 0.95  & \textbf{I.P.\ SNIP-it (ours)}   & \result{vgg19.ip.snip.0.95} & this paper \\
VGG-19 & 0.95  & \textbf{BaCP SNIP-it (ours)}    & \result{vgg19.bacp.snip.0.95} & this paper \\
\midrule
\multicolumn{5}{l}{\textit{ResNet-34 / CIFAR-10 --- extreme sparsity, dynamic methods only (context)}} \\
ResNet-34 & 0.99  & RigL (dynamic)   & $92.92 \pm 0.18$ & \citep{li2025east} \\
ResNet-34 & 0.99  & EAST (dynamic)   & $93.51 \pm 0.13$ & \citep{li2025east} \\
ResNet-34 & 0.999 & RigL (dynamic)   & $85.71 \pm 0.23$ & \citep{li2025east} \\
ResNet-34 & 0.999 & EAST (dynamic)   & $86.99 \pm 0.32$ & \citep{li2025east} \\
ResNet-34 & 0.999 & SET (dynamic)    & $82.70 \pm 0.91$ & \citep{li2025east} \\
ResNet-34 & 0.99  & \textbf{I.P.\ magnitude (ours, static)} & \result{resnet34.ip.magnitude.0.99} & this paper \\
ResNet-34 & 0.99  & \textbf{BaCP magnitude (ours, static)}  & \result{resnet34.bacp.magnitude.0.99} & this paper \\
ResNet-34 & 0.999 & \textbf{I.P.\ magnitude (ours, static)} & \result{resnet34.ip.magnitude.0.999} & this paper \\
ResNet-34 & 0.999 & \textbf{BaCP magnitude (ours, static)}  & \result{resnet34.bacp.magnitude.0.999} & this paper \\
\bottomrule
\end{tabular}
\end{table}

The GMP row is the one to attend to. Gradual magnitude pruning is our
\texttt{magnitude} criterion under a different name --- \citet{zhu2018prune}'s
cubic schedule applied to $|W|$ --- so the ResNet-50 / 95\% GMP entry is a direct
match for our I.P.\ magnitude cell in method, model, dataset and sparsity, and
differs from it only in budget. It is the single most informative published number
for this paper: it tells a reader whether the baseline BaCP is measured against is
a competently trained one.
% TODO: confirm direction once the sweep lands -- state the actual gap between our I.P. magnitude cell and published GMP, and attribute it to budget only if the size is consistent with the budget difference.

The ResNet-34 block is context, not comparison, and we label it as such: EAST and
RigL are \emph{dynamic} sparse training methods that reallocate their support
during training, whereas both of our arms are static, and no published static
ResNet-34 / CIFAR-10 result exists at any sparsity to compare against. We include
the block because it is the only published evidence of what 99.9\% costs a ResNet
on this dataset, and a reader is entitled to that scale.
% TODO: confirm direction once the sweep lands -- do not claim parity with EAST/RigL at 0.999; if our static numbers sit well below, say so plainly.

Two further gaps we state rather than paper over. VGG-11 has \textbf{no published
comparator at any sparsity}, so its entire column in Table~\ref{tab:main} is
reported standalone. And no published static comparator exists at 0.97, 0.99 or
0.999 for any of our four models, so for three of our four sparsity levels the
I.P.\ baseline in Table~\ref{tab:main} is the only reference point --- which is
exactly why that baseline was run under a matched budget rather than quoted.
